\section{Distribution Recursions}
\label{sec:distribution-recursions}

\setcounter{exercise}{88}
\begin{exercise}
\label{ex:distribution-recursions-89}
Let $I_1,\ldots,I_n$ be independent random variables, each of which is equally likely to be either $0$ or $1$. A well-known nonparametric statistical test, called the signed rank test, is concerned with determining $P_n(k)$ defined by
\[
P_n(k)=P\left\{\sum_{j=1}^n jI_j\le k\right\}.
\]
Justify the following formula:
\[
P_n(k)=\frac12 P_{n-1}(k)+\frac12 P_{n-1}(k-n).
\]
\end{exercise}

\solution{
Let
\[
S_n=\sum_{j=1}^n jI_j.
\]
Then
\[
S_n=S_{n-1}+nI_n,
\]
where
\[
S_{n-1}=\sum_{j=1}^{n-1}jI_j.
\]

Since $I_n$ is equally likely to be $0$ or $1$, we condition on the value of $I_n$:
\[
P_n(k)
=
P(S_n\le k).
\]
Thus,
\[
P_n(k)
=
P(S_n\le k\mid I_n=0)P(I_n=0)
+
P(S_n\le k\mid I_n=1)P(I_n=1).
\]

Since
\[
P(I_n=0)=P(I_n=1)=\frac12,
\]
we get
\[
P_n(k)
=
\frac12 P(S_{n-1}\le k)
+
\frac12 P(S_{n-1}+n\le k).
\]
The first term is
\[
P(S_{n-1}\le k)=P_{n-1}(k),
\]
and the second term is
\[
P(S_{n-1}+n\le k)=P(S_{n-1}\le k-n)=P_{n-1}(k-n).
\]
Therefore,
\[
P_n(k)=\frac12 P_{n-1}(k)+\frac12 P_{n-1}(k-n).
\]
}
